% !TEX program = pdflatex
% =============================================================================
% Corporate Bankruptcy Prediction: Econometric Benchmarks vs. Machine Learning
% Sections 1-3 initial draft
% Compiles as-is in Overleaf with no extra packages required.
% =============================================================================

\documentclass[11pt,a4paper]{article}

% -------- packages --------
\usepackage[utf8]{inputenc}
\usepackage[T1]{fontenc}
\usepackage{lmodern}
\usepackage[a4paper,margin=2.5cm]{geometry}
\usepackage{amsmath,amssymb,amsthm}
\usepackage{booktabs}
\usepackage{graphicx}
\usepackage{array}
\usepackage{tabularx}
\usepackage{multirow}
\usepackage{caption}
\usepackage[round]{natbib}
\usepackage{setspace}
\usepackage{csquotes}
\usepackage{hyperref}
\hypersetup{
  colorlinks=true,
  linkcolor=black,
  citecolor=black,
  urlcolor=blue,
}

\onehalfspacing
\setlength{\parskip}{4pt}
\setlength{\parindent}{16pt}

% -------- title block --------
\title{\textbf{Modeling the Probability of Corporate Default Using Econometric and Machine Learning Methods}}
\author{Completed by: Ilias Tagirov \\ \small
National Research University Higher School of Economics \\ \small
International Institute of Economics and Finance \\ \small
Scientific Advisor: Nadezhda Kanygina}
\date{\today}

\begin{document}
\maketitle

\begin{abstract}
\noindent
This paper compares classical accounting-based bankruptcy prediction models
--- \cite{altman1968}, \cite{ohlson1980}, and \cite{zmijewski1984} --- with
modern machine-learning (ML) classifiers on the publicly available Polish
Companies Bankruptcy dataset of \cite{zieba2016}. Unlike prior benchmark
comparisons in this literature, which often vary imbalance-handling, feature
engineering, and evaluation protocols across models, we fix a single protocol
across all specifications: stratified five-fold cross-validation, class-weighted
loss functions (no undersampling of the healthy class), median imputation fit
inside each training fold, and identical discrimination, calibration, and
economic-cost metrics for every model and every horizon. We re-estimate the
three canonical econometric benchmarks on the Polish data, compare them to
penalized logistic regression and gradient-boosted trees over five prediction
horizons (1 to 5 years ahead), and decompose the sources of any ML advantage
into discrimination, calibration, and false-alarm suppression. The resulting
comparison is intended to be more credible as a basis for practitioner and
academic decisions than existing benchmark claims in which protocol differences
confound the model-vs-model comparison.
\end{abstract}

\vspace{1em}

% =============================================================================
\section{Introduction}
% =============================================================================

Predicting corporate bankruptcy is one of the oldest empirical tasks in
finance and one of the most consequential. A firm's default imposes losses on
its creditors, suppliers, employees, and shareholders, and in the aggregate
bankruptcies transmit financial shocks across the real economy. The 2007--2009
financial crisis made these concerns acutely visible: public concern with
distress risk, regulatory interest in early-warning indicators, and private
interest from banks and insurers in internal-rating models all intensified
sharply after the crisis, producing an empirical literature that now spans
more than half a century of model development
\citep{altman1968, ohlson1980, zmijewski1984, shumway2001, campbell2008,
barboza2017, zieba2016, mai2019}.

\subsection{The state of the literature}

The literature divides into two broad methodological camps. The first,
\textbf{econometric}, is grounded in financial-statement analysis and
parametric statistical models. Its canonical entries are
\cite{altman1968}, who applied multiple discriminant analysis to five
accounting ratios to classify 33 bankrupt versus 33 non-bankrupt U.S.
manufacturing firms; \cite{ohlson1980}, who replaced the discriminant analysis
with a logistic regression (logit) on nine variables and a much larger sample
(105 bankrupt vs.\ 2,058 non-bankrupt U.S.\ industrial firms, 1970--1976); and
\cite{zmijewski1984}, who used a probit regression on three ratios in a
sample of 40 bankrupt and 800 non-bankrupt firms. \cite{shumway2001} later
generalized the single-period logit into a discrete-time hazard model,
allowing predictors to change over time and handling censoring explicitly.
These models share a common set of virtues: they are interpretable, their
coefficients have direct economic meaning (a negative coefficient on
retained-earnings-to-assets, for example, means cumulative profitability
reduces bankruptcy risk), and their standard errors admit the usual
hypothesis-testing toolkit.

The second camp is \textbf{machine learning} (ML). A stream of research
beginning in the 1990s and accelerating after the 2007--2009 crisis has
applied progressively more flexible classifiers to bankruptcy prediction:
neural networks \citep{altman1994}, support vector machines
\citep{shin2005, min2005}, random forests and ensemble methods
\citep{barboza2017}, gradient-boosted decision trees
\citep{zieba2016, carmona2019}, and more recently deep sequence models such as
LSTMs \citep{mai2019, pellegrino2024}. The empirical consensus of this camp
is that ML methods attain higher out-of-sample discrimination, typically
measured by the area under the ROC curve (AUC), than the classical
econometric benchmarks they are compared against.

\subsection{A methodological concern with existing comparisons}

The ``ML is better than logit'' finding is now widespread in the bankruptcy
prediction literature. On closer inspection, however, several methodological
choices typically differ between the ML and econometric models in the
same study, making the model-vs-model comparison confounded. Three confounds
are particularly common.

First, \textbf{imbalance handling}. Bankruptcy is a rare event: annual
default rates for listed firms are typically below 5\% and often below 1\%
\citep{shumway2001}. Many ML studies aggressively resample the training data
--- either by undersampling the healthy class, as in \cite{pellegrino2022}, or
by synthetic oversampling via SMOTE \citep{chawla2002} --- while the
econometric benchmark is estimated on the raw imbalanced sample. Resampling
has well-known consequences for calibration: a classifier trained on a
50/50 balanced sample will, by construction, produce probability estimates
that are systematically too high for new data whose true base rate is 4\%
\citep{king2001, saerens2002}. An ML model that has been recalibrated, even
implicitly, towards a different base rate will score higher on discrimination
metrics that reward ranking but may fail calibration metrics that reward
correct absolute probability.

Second, \textbf{feature set}. ML models are often given access to many more
predictors than the accounting ratios on which Altman or Ohlson was originally
estimated. An ML model with 64 features and 500 hyperparameter configurations
that beats a 5-variable Altman logit provides, strictly speaking, evidence
about feature engineering and hyperparameter search rather than about the
inherent superiority of the ML algorithm.

Third, \textbf{evaluation metric}. Discrimination metrics such as the AUC
and the precision-recall AUC (PR-AUC) measure how well a model \emph{ranks}
firms by bankruptcy risk; calibration metrics such as the Brier score and
the expected calibration error measure how accurate a model's
\emph{probability estimates} are in absolute terms; and economic cost metrics
combine misclassifications with a loss function that reflects the real-world
cost asymmetry between false positives and false negatives. These metrics
capture genuinely different properties of a classifier, and it is possible
(and common) for one model to beat another on AUC while losing on calibration.
Reporting a single metric, typically AUC, hides this trade-off.

\subsection{This paper}

We re-examine the ML-vs-econometric question on a single widely used public
dataset, the Polish Companies Bankruptcy dataset of \cite{zieba2016}, under
one fixed evaluation protocol applied identically to every model. Concretely,
our contributions are as follows.

\begin{enumerate}
\item We re-estimate the three canonical econometric benchmarks
(\citealt{altman1968}, \citealt{ohlson1980}, \citealt{zmijewski1984}) on the
Polish data using their original variable sets where those variables are
available, and we document explicitly which original variables are and are
not recoverable.

\item We compare these benchmarks with a penalized logistic regression
(LASSO), a random forest, and an extreme gradient-boosted tree model
(XGBoost), under one fixed protocol: stratified five-fold cross-validation,
class-weighted loss (no undersampling, no SMOTE in the primary specification),
median imputation of missing ratios fit on the training fold only, and an
identical feature set for every model of the same complexity class.

\item We evaluate every model on a standard battery of metrics spanning
three conceptually distinct properties: discrimination (AUC, PR-AUC, KS
statistic), calibration (Brier score, reliability diagrams, expected
calibration error, post-hoc isotonic recalibration), and an economic cost
metric based on the loss-ratio assumptions used in the credit-risk
literature \citep{altman2007}.

\item We exploit the panel structure of the Polish dataset --- which contains
five parallel files, one per prediction horizon from 1 year to 5 years ahead
--- to test whether the ML advantage persists as the horizon lengthens, or
whether (as several classical studies have suggested)
\citep{mossman1998, duan2012} all models converge towards uninformative
predictions at long horizons.

\item We quantify the sources of any residual ML advantage, distinguishing
explicitly whether ML wins by detecting more bankrupt firms (higher
bankruptcy-class recall) or by reducing false alarms among healthy firms
(higher bankruptcy-class precision at equal recall).

\item We document the informative-missingness structure of the Polish data
--- in short, distressed firms systematically fail to report certain
line items --- and show that treating missingness as a feature in its own
right (via binary indicator variables) is both a defensible modeling choice
and an empirically important one.
\end{enumerate}

The paper is organized as follows. Section~\ref{sec:lit} surveys the
accounting-based, statistical, and ML literatures on bankruptcy prediction,
and critiques their comparative methodology. Section~\ref{sec:data} describes
the Polish Companies Bankruptcy dataset. Section~\ref{sec:methodology}
specifies the six models and the unified evaluation protocol. Section 5
reports the results, Section 6 interprets them, Section 7 presents robustness
checks, and Section 8 concludes.

% =============================================================================
\section{Literature Review}\label{sec:lit}
% =============================================================================

The bankruptcy-prediction literature is vast; we organize our review into
three connected blocks. First, the classical accounting-based scoring models,
which established the intellectual framework for the field and whose
re-estimated versions are the econometric benchmarks of this paper. Second,
the more recent ML contributions, which have established the ``ML beats
econometrics'' empirical stylized fact we propose to re-examine. Third, the
methodological literature on imbalanced classification in finance, which
motivates our unified evaluation protocol.

\subsection{Classical accounting-based scoring models}

\paragraph{Altman (1968): multiple discriminant analysis.}
\cite{altman1968} is widely regarded as the first rigorous multivariate
statistical model of corporate bankruptcy. On a sample of 66 publicly traded
U.S.\ manufacturing firms --- 33 that had filed under Chapter~X of the 1938
Bankruptcy Act between 1946 and 1965, matched to 33 non-bankrupt firms by
industry and size --- Altman applied \textit{multiple discriminant analysis}
(MDA), a technique introduced by \cite{fisher1936}, to select five
accounting ratios from an initial pool of 22. MDA estimates the linear
combination of predictors that best separates two (or more) class means
relative to within-class variance, subject to the assumption that the
predictors are jointly normally distributed and have equal covariance
matrices across classes. The resulting \textit{Z-score} is
\begin{equation}
Z = 1.2 \cdot X_1 + 1.4 \cdot X_2 + 3.3 \cdot X_3 + 0.6 \cdot X_4 + 1.0 \cdot X_5
\label{eq:altman}
\end{equation}
where $X_1$ is working capital over total assets (a liquidity measure), $X_2$
is retained earnings over total assets (a cumulative profitability and
firm-age measure), $X_3$ is earnings before interest and taxes (EBIT) over
total assets (an operating profitability measure), $X_4$ is the market value
of equity over the book value of total liabilities (a leverage-and-solvency
measure), and $X_5$ is sales over total assets (an activity measure).

Firms were classified as likely bankrupt if $Z < 1.81$ and as likely solvent
if $Z > 2.99$; values between were labelled a ``zone of ignorance.''
\cite{altman1968} reported 95\% classification accuracy one year before
bankruptcy and 72\% two years before, on his matched-pair sample.
Altman's framework has been extended many times --- \cite{altman2000} for
private firms (the Z'-score), for non-manufacturing firms (the Z''-score), and
for emerging markets --- and remains a standard professional tool in credit
analysis. Its key methodological limitations are the assumptions of joint
normality and equal covariances (both violated by accounting ratios), the
bias introduced by matched-sample estimation when the population base rate is
far from 50\%, and the requirement that market value be observable (excluding
private firms in the original formulation).

\paragraph{Ohlson (1980): logistic regression.}
\cite{ohlson1980} relaxed the distributional assumptions of MDA by replacing
it with a logistic regression. The logit model posits
\begin{equation}
\Pr(y_i = 1 \mid \mathbf{x}_i) = \frac{1}{1 + \exp(-\mathbf{x}_i^\top \boldsymbol{\beta})},
\label{eq:logit}
\end{equation}
where $y_i = 1$ if firm $i$ goes bankrupt and $\mathbf{x}_i$ is a vector of
predictors. The coefficients $\boldsymbol{\beta}$ are estimated by maximum
likelihood, with no normality requirement on the predictors and a direct
probability interpretation of the output.

Ohlson estimated the model on 105 bankrupt and 2,058 non-bankrupt U.S.\
industrial firms drawn from the Compustat database over the fiscal years
1970--1976, using nine predictors that he called the ``O-score'' components:
the log of total assets scaled by the GNP deflator (SIZE); total liabilities
over total assets (TLTA); working capital over total assets (WCTA); current
liabilities over current assets (CLCA); a dummy equal to one if total
liabilities exceed total assets (OENEG); net income over total assets (NITA);
funds from operations over total liabilities (FUTL); a dummy equal to one if
net income was negative in each of the two preceding years (INTWO); and a
normalized measure of change in net income (CHIN). The three innovations
relative to Altman --- larger and more representative sample, relaxation of
the normality and equal-covariance assumptions, and direct probability output
--- have made Ohlson's O-score the most widely used accounting-based
bankruptcy model in academic research. It is the natural econometric
benchmark for any modern bankruptcy-prediction study.

\paragraph{Zmijewski (1984): probit regression.}
\cite{zmijewski1984} proposed a parsimonious three-variable probit model.
The probit differs from the logit only in its link function: instead of the
logistic CDF in Equation~\ref{eq:logit}, it uses the standard normal CDF,
\begin{equation}
\Pr(y_i = 1 \mid \mathbf{x}_i) = \Phi\!\left( \mathbf{x}_i^\top \boldsymbol{\beta} \right).
\label{eq:probit}
\end{equation}
In most applied settings the logit and probit produce extremely similar
predictions, and the choice between them is largely a matter of convention;
we keep Zmijewski as a probit for fidelity to the original paper. Zmijewski's
three predictors are net income over total assets, total liabilities over
total assets, and current assets over current liabilities. He estimated the
model on 40 bankrupt and 800 non-bankrupt U.S.\ firms over 1972--1978. Besides
its parsimony, the Zmijewski model is notable for explicitly addressing the
\textit{choice-based sampling bias} that arises when bankrupt firms are
oversampled relative to their population prevalence: he showed that the slope
coefficients are consistent under choice-based sampling even when the
intercept is not, which remains a standard result in the literature.

\paragraph{Shumway (2001): discrete-time hazard.}
\cite{shumway2001} argued that the earlier cross-sectional models are
statistically inconsistent for the underlying problem: bankruptcy is a
time-to-event process, firms are observed for varying numbers of years, and
predictors change over time. He proposed a discrete-time hazard model that
can be estimated by pooled logit with an auxiliary time variable. On a panel
of U.S.\ firms 1962--1992 he found that roughly half the accounting ratios
that appeared significant in earlier cross-sectional models lost
significance once the hazard structure was imposed, and that market variables
(equity market capitalization, past returns, idiosyncratic return volatility)
were strong and robust additional predictors. We note this model here as the
conceptually correct framework for panel bankruptcy data, but we cannot
estimate it on the Polish Companies dataset because that dataset does not
contain firm identifiers that would allow a panel to be constructed
(Section~\ref{sec:data}).

\subsection{Machine-learning bankruptcy models}

The application of ML to bankruptcy prediction began in the late 1980s
with early neural networks \citep{odom1990} and has proliferated steadily
since. We highlight the threads most relevant to this paper.

\cite{barboza2017} conducted a broad comparison of ML methods (support
vector machines, random forests, bagging, boosting, and neural networks)
against discriminant analysis and logistic regression on a large sample of
U.S.\ firms. They reported that ML ensembles achieved roughly ten percentage
points higher accuracy than the classical models. Their study is often cited
as strong evidence for ML superiority, but the authors did not publicly
release their data or the specific hyperparameter settings of their models,
which precludes exact replication.

\cite{zieba2016}, whose dataset we use in this paper, applied the then-new
XGBoost gradient-boosting algorithm to Polish firms over 2000--2013. They
reported AUCs in the 0.93--0.96 range on the five forecasting horizons
(1 to 5 years ahead), materially higher than the linear logistic regression
baselines they compared against. They attributed the advantage to the
ability of tree-based ensembles to handle the highly skewed distributions of
accounting ratios without normalization, to capture nonlinearities, and to
their low sensitivity to class imbalance.

\cite{pellegrino2022} applied deep neural networks and tree ensembles to
a dataset of 8,262 U.S.\ public companies over 1999--2018, with five
prediction horizons. Their best neural network achieved a test-set AUC of
0.85 and was judged to outperform their random-forest baseline (mean AUC
0.748). However, their training protocol imposed random undersampling of
the healthy class to achieve a 50/50 training balance, separately for every
one of their 100 repeated training runs, before evaluating on the original
imbalanced validation and test sets. As \cite{king2001} have emphasized,
this protocol systematically biases the predicted probabilities upward and
complicates direct comparison to models trained on the natural class
prevalence. We treat this paper as an important reference point and a
cautionary example of the confounding between imbalance handling and model
class that motivates our unified protocol.

\cite{mai2019} apply deep-learning techniques to textual disclosures from
U.S.\ 10-K filings, and show that text adds incremental predictive power
beyond accounting variables. This line of work expands the input space
beyond structured financial data; our paper deliberately restricts itself to
structured accounting ratios to isolate the model-class comparison from
confounding gains due to new data sources.

\subsection{Imbalanced classification and evaluation}

Bankruptcy prediction is a textbook imbalanced classification problem, and a
substantial methodological literature warns against naïve handling of class
imbalance. \cite{saito2015} show that the PR-AUC is preferable to the ROC-AUC
for evaluating classifiers on heavily imbalanced data, because the ROC curve
can look deceptively good when the negative class dominates, while the PR
curve exposes the difficulty of precision on the minority class.
\cite{king2001} document the systematic probability inflation that results
from choice-based sampling or undersampling, and provide a simple post-hoc
correction for the intercept; this correction is widely ignored in practice.
\cite{niculescu2005} demonstrate that tree-based ensembles (random forests,
boosted trees), while strong on discrimination, are often poorly calibrated
out of the box, and that post-hoc isotonic or Platt recalibration
substantially improves their Brier scores. \cite{altman2007} discuss
asymmetric misclassification costs in credit risk: a false negative (a
predicted-healthy firm that actually defaults) is typically 20 to 50 times
more costly than a false positive (a predicted-bankrupt firm that in fact
survives), so evaluation under a fixed 0/1 threshold masks economically
important trade-offs.

These three strands --- imbalance-aware metrics, sampling-induced calibration
bias, and asymmetric costs --- motivate the specific evaluation protocol we
adopt in Section~\ref{sec:methodology}. In brief: we report AUC \textit{and}
PR-AUC \textit{and} Brier score \textit{and} a cost-weighted misclassification
metric for every model, we estimate every model on the natural class
prevalence with a class-weighted loss rather than through resampling, and we
apply post-hoc isotonic recalibration as a separate reported step rather than
as a hidden pre-processing choice.

% =============================================================================
\section{Data}\label{sec:data}
% =============================================================================

\subsection{Source and provenance}

We use the \textit{Polish Companies Bankruptcy} dataset compiled by
\cite{zieba2016} and released publicly through the UCI Machine Learning
Repository (dataset ID~365).\footnote{\url{https://doi.org/10.24432/C5F600};
CC~BY~4.0 licence.} The underlying data were collected from the Emerging
Markets Information Service (EMIS), a commercial database of emerging-market
corporate filings. Non-bankrupt firms are observed for fiscal years
2007--2013; bankrupt firms are observed for fiscal years 2000--2012 prior to
their bankruptcy filings.

The dataset contains 64 pre-computed financial ratios (liquidity,
profitability, leverage, activity, and efficiency indicators) for each
firm-year observation. Crucially for this paper, \textit{the ratios are
pre-computed by the original authors from the underlying financial
statements}: we do not observe the raw balance-sheet and income-statement
items, and we therefore do not face the unit-consistency problems that can
arise when ratios are constructed on the fly from possibly differently scaled
raw fields. This is a meaningful practical advantage over other public
bankruptcy datasets that expose raw accounting items only.

\subsection{Horizon structure and internal renaming}
\label{sec:horizons}

The dataset is distributed as five separate files, each representing a
different prediction-horizon experiment. The original UCI naming is
counterintuitive: the file named \texttt{1year.arff} contains financial
ratios from the \textit{earliest} available fiscal year of the forecasting
window, with a binary label indicating whether the firm went bankrupt
\textit{within the subsequent five years}. Correspondingly
\texttt{5year.arff} contains ratios from the latest available fiscal year,
with a label indicating whether the firm went bankrupt \textit{within the
subsequent one year}. To avoid this confusion throughout the paper, we
relabel the files internally by the length of the prediction horizon, where
a shorter horizon means the financial statements are closer in time to the
observed bankruptcy outcome:

\begin{center}
\begin{tabular}{lll}
\toprule
UCI file & Our label & Prediction horizon \\
\midrule
\texttt{5year.arff} & H1 & 1 year ahead  \\
\texttt{4year.arff} & H2 & 2 years ahead  \\
\texttt{3year.arff} & H3 & 3 years ahead  \\
\texttt{2year.arff} & H4 & 4 years ahead  \\
\texttt{1year.arff} & H5 & 5 years ahead  \\
\bottomrule
\end{tabular}
\end{center}

\noindent Under this convention H1 is the shortest (and highest-signal)
prediction task and H5 is the longest (and lowest-signal) one. H1 is our
primary analysis horizon; H2--H5 are used in the horizon-degradation tests
of Section~5.

\subsection{Sample sizes and class balance}

Table~\ref{tab:horizons} summarizes the sample sizes, bankruptcy counts, and
base rates by horizon. Three patterns are visible. First, the bankruptcy
rate is low and falls substantially from 6.9\% at H1 to 3.9\% at H5. This
reflects a combination of the shorter-horizon samples being selected for
firms whose bankruptcy is imminent (and therefore more likely to appear in
the bankrupt subset) and the longer-horizon samples including more firms
that will eventually fail but are not yet in the distress phase. Second, all
five horizons have class imbalance between roughly 14:1 (H5) and 14:1 (H1)
healthy:bankrupt, which is severe for an ML classification problem. Third,
the sample sizes are smaller than would be ideal: H1 has only 410 bankrupt
firms across 5,910 observations, and H5 has 271 across 7,027. This
sample-size limitation must be remembered when evaluating any complex
ML classifier's claim to ``capture nonlinear interactions:'' the effective
training signal is driven by the positive class, whose size is in the
low hundreds.

\begin{table}[ht]
\centering
\caption{Polish Companies Bankruptcy: sample sizes and class balance by
prediction horizon. Bankruptcy-rate-among-complete-rows and
bankruptcy-rate-among-rows-with-any-missing-value are reported to document
the informative-missingness pattern discussed in Section~\ref{sec:missing}.}
\label{tab:horizons}
\begin{tabular}{lrrrrrr}
\toprule
 & Obs. & Bankrupt & Rate (\%) & Complete rows & BR complete (\%) & BR incomplete (\%) \\
\midrule
H1 (1 yr ahead) & 5{,}910  & 410 & 6.94 & 3{,}031 & 3.37 & 10.70 \\
H2 (2 yr ahead) & 9{,}792  & 515 & 5.26 & 4{,}769 & 2.52 & 7.86 \\
H3 (3 yr ahead) & 10{,}503 & 495 & 4.71 & 4{,}885 & 2.19 & 6.91 \\
H4 (4 yr ahead) & 10{,}173 & 400 & 3.93 & 4{,}088 & 1.79 & 5.37 \\
H5 (5 yr ahead) & 7{,}027  & 271 & 3.86 & 3{,}194 & 0.94 & 6.29 \\
\bottomrule
\end{tabular}
\end{table}

\subsection{The 64 financial ratios}
\label{sec:ratios}

Each firm-year observation in the dataset comprises 64 ratios derived from
the firm's annual financial statements. The ratios were selected by
\cite{zieba2016} to span the standard categories used in the accounting-based
bankruptcy literature: liquidity, profitability, leverage, activity, and
efficiency. In our analysis we retain all 64 ratios and rename them to
interpretable labels reflecting their economic meaning (e.g.\ we refer to
\textit{Attr1} as \texttt{net\_profit\_TA}, to \textit{Attr7} as
\texttt{EBIT\_TA}, etc.).

Importantly, the 64 ratios include \textit{all five} components of the
\cite{altman1968} Z-score:
\begin{itemize}
\item \textbf{Working capital / Total assets} --- Altman's $X_1$, our
\texttt{WC\_TA} (Attr3).
\item \textbf{Retained earnings / Total assets} --- Altman's $X_2$, our
\texttt{RE\_TA} (Attr6).
\item \textbf{EBIT / Total assets} --- Altman's $X_3$, our \texttt{EBIT\_TA}
(Attr7).
\item \textbf{Book value of equity / Total liabilities} --- Altman's $X_4$
(adapted for non-listed firms), our \texttt{BE\_TL} (Attr8).
\item \textbf{Sales / Total assets} --- Altman's $X_5$, our \texttt{sales\_TA}
(Attr9).
\end{itemize}

\noindent
Seven of the nine components of \cite{ohlson1980}'s O-score are directly
available (SIZE, TLTA, WCTA, CLCA, NITA, FUTL, and implicitly OENEG); the
two that are not (INTWO and CHIN) require consecutive firm-year observations
with firm identifiers, which this dataset does not supply.\footnote{The
dataset is anonymized at the firm level: each row is a firm-year observation,
but different rows cannot be matched to the same firm. This prevents the
construction of within-firm lag features and, as noted in
Section~\ref{sec:lit}, precludes estimation of a \cite{shumway2001}
discrete-time hazard model on these data.} All three components of the
\cite{zmijewski1984} X-score are directly available (NITA, TLTA, CACL).
Our re-estimated econometric specifications are therefore faithful
replications of Altman and Zmijewski, and a seven-of-nine replication of
Ohlson.

\subsection{Missing values and the informative-missingness structure}
\label{sec:missing}

Unlike some widely used public datasets (for instance the recent U.S.\
benchmark of \citealt{pellegrino2022}, in which firm-years with any missing
ratio are dropped prior to release), the Polish Companies data retain rows
with missing values. The missingness rate per feature, computed across all
five horizons, is low for most features --- 49 of 64 features have
essentially no missingness --- but concentrated in a small set of ratios
whose denominators involve components that can legitimately be zero or
unobserved. The most heavily missing feature is \texttt{CA\_inv\_LTL}
(Attr37: current assets minus inventories, divided by long-term liabilities),
which is missing in 43\% of H1 rows, simply because a firm with no long-term
debt has an undefined ratio.

Crucially, this missingness is \textit{informative}: it is not random with
respect to the bankruptcy outcome. Table~\ref{tab:horizons} contrasts the
bankruptcy rate among firms with fully complete observations against that
among firms with at least one missing ratio. In every horizon the incomplete
group has a higher bankruptcy rate than the complete group, and the gap is
large:

\begin{itemize}
\item At H1 the incomplete-row bankruptcy rate (10.7\%) is 3.2 times the
complete-row rate (3.4\%).
\item At H5 the gap is most extreme: incomplete-row bankruptcy rate of 6.3\%
against a complete-row rate of 0.9\%, a factor of 6.7.
\end{itemize}

\noindent This pattern is consistent with the accounting intuition that
firms in financial distress are more likely to have degraded reporting: they
may lose long-term debt (driving \texttt{CA\_inv\_LTL} to undefined), they
may restructure asset categories in ways that blank certain items, or their
reporting discipline simply declines. The missingness is therefore a signal
about distress in its own right, and treating it with naïve imputation ---
replacing NaNs with the sample median and proceeding as if the value were
observed --- would discard real information.

We address this in two ways, both motivated by the credit-risk and
missing-data literatures. First, for every feature that is missing in at
least 1\% of rows in any horizon, we add a binary indicator column
\texttt{miss\_<feature>} that takes the value one when the original feature
is missing and zero otherwise. Eleven such indicators are constructed.
These indicators are deterministic functions of the raw-data NaN pattern ---
they involve no fitting and therefore no risk of data leakage --- and they
are included as first-class features in every model, alongside the 64 raw
ratios. Second, we impute the remaining missing values in the raw ratios
with the \textit{training-fold median}, fit inside each cross-validation
fold (Section~\ref{sec:methodology}). Median imputation is simple, robust to
outliers, and standard in credit-risk work; the imputed value is not itself
the signal --- the indicator is. As a robustness check
(Section~7) we repeat the main analysis on the complete-case subsample only,
so that the reader can assess whether our results depend on the imputation.

\subsection{Panel structure and its limits}

As noted in Section~\ref{sec:ratios}, the dataset does not contain firm
identifiers that would link observations across years for the same firm.
Each of the five horizon files is therefore a cross-sectional dataset in its
own right; there is no panel time dimension within a file, and no way to
follow a firm from one file to the next. This has three consequences for our
empirical design. First, we cannot estimate a discrete-time hazard model in
the sense of \cite{shumway2001} --- the relevant transition probabilities are
not identified without the panel structure. Second, we cannot construct
within-firm lagged features, which rules out the two Ohlson variables
INTWO and CHIN. Third, the five horizon files are plausibly
\textit{overlapping}: the same firm may appear in H1, H2, and H3 under
different fiscal-year observations, potentially with different labels as its
distress evolves. We are not able to deduplicate across horizons, and this
is a limitation of the analysis.

We rely instead on a \textit{stratified five-fold cross-validation}
protocol within each horizon file: each firm-year observation is assigned
to one of five folds with stratification on the bankruptcy label so that the
class balance is approximately preserved in every fold. Folds are stratified
(not random) because with only 271--515 bankrupt firms in each horizon, a
random split could produce folds with substantially different minority-class
prevalence, which would add noise to the fold-level metric estimates.

% =============================================================================
% PLACEHOLDER: Section 3.6 descriptive statistics table
% To be filled in after Cell 3 runs.
% =============================================================================
\subsection{Descriptive statistics of the Altman ratios}
\label{sec:descriptives}

Table~\ref{tab:descriptives} reports, for each of the five Altman ratios,
the mean and standard deviation among bankrupt and non-bankrupt firms
separately at horizon H1, together with a Welch two-sample $t$-statistic
testing the null of equal means across the two groups. (Welch's $t$-test is
used because the within-group variances differ markedly, violating the
equal-variance assumption of the standard $t$-test; the Welch version uses
a modified degrees-of-freedom formula to handle unequal variances.)

\begin{table}[ht]
\centering
\caption{Altman ratio means by bankruptcy status, horizon H1. Welch's
$t$-statistic tests equality of means between the bankrupt and non-bankrupt
groups; a large absolute $t$ indicates that the ratio differs materially
between groups and is therefore a candidate for inclusion in a bankruptcy
model.}
\label{tab:descriptives}
\begin{tabular}{lcccccc}
\toprule
Ratio & Mean (non-bankrupt) & SD & Mean (bankrupt) & SD & $t$-stat & $p$ \\
\midrule
WC\_TA       & TBD & TBD & TBD & TBD & TBD & TBD \\
RE\_TA       & TBD & TBD & TBD & TBD & TBD & TBD \\
EBIT\_TA     & TBD & TBD & TBD & TBD & TBD & TBD \\
BE\_TL       & TBD & TBD & TBD & TBD & TBD & TBD \\
sales\_TA    & TBD & TBD & TBD & TBD & TBD & TBD \\
\bottomrule
\end{tabular}
\caption*{\footnotesize Values to be filled in after Cell~3 runs.
$t$-statistics with absolute value above roughly 2 indicate that the ratio
differs significantly between the two groups at a 5\% level.}
\end{table}

Interpretation --- to be written once numbers are available:
\begin{itemize}
\item \textit{Expected sign} of each ratio's difference between bankrupt and
non-bankrupt firms. WC\_TA, RE\_TA, EBIT\_TA, and BE\_TL should be smaller
for bankrupt firms. sales\_TA is ambiguous: low sales relative to assets
indicate inefficient asset utilization but high sales relative to assets can
also reflect distressed asset sales.
\item \textit{Which ratios discriminate most strongly?} If EBIT\_TA and
RE\_TA show the largest $|t|$ and WC\_TA and sales\_TA the smallest, this
is consistent with \cite{altman1968}'s own finding that cumulative and
operating profitability dominate the Z-score.
\item \textit{Implications for the downstream Altman logit coefficients.}
The signs observed in this descriptive table should match the signs of the
re-estimated logit coefficients reported in Section~5; any discrepancies
deserve discussion.
\end{itemize}

% =============================================================================
% Bibliography
% =============================================================================

\bibliographystyle{abbrvnat}
\begin{thebibliography}{99}

\bibitem[Altman(1968)]{altman1968}
Altman, E.\ I.\ (1968).
\newblock Financial ratios, discriminant analysis and the prediction of
corporate bankruptcy.
\newblock \textit{Journal of Finance}, 23(4), 589--609.

\bibitem[Altman et~al.(1994)]{altman1994}
Altman, E.\ I., Marco, G., and Varetto, F.\ (1994).
\newblock Corporate distress diagnosis: comparisons using linear discriminant
analysis and neural networks (the Italian experience).
\newblock \textit{Journal of Banking \& Finance}, 18(3), 505--529.

\bibitem[Altman(2000)]{altman2000}
Altman, E.\ I.\ (2000).
\newblock Predicting financial distress of companies: revisiting the Z-score
and ZETA models.
\newblock Working paper, Stern School of Business, New York University.

\bibitem[Altman and Sabato(2007)]{altman2007}
Altman, E.\ I.\ and Sabato, G.\ (2007).
\newblock Modelling credit risk for SMEs: evidence from the US market.
\newblock \textit{Abacus}, 43(3), 332--357.

\bibitem[Barboza et~al.(2017)]{barboza2017}
Barboza, F., Kimura, H., and Altman, E.\ (2017).
\newblock Machine learning models and bankruptcy prediction.
\newblock \textit{Expert Systems with Applications}, 83, 405--417.

\bibitem[Campbell et~al.(2008)]{campbell2008}
Campbell, J.\ Y., Hilscher, J., and Szilagyi, J.\ (2008).
\newblock In search of distress risk.
\newblock \textit{Journal of Finance}, 63(6), 2899--2939.

\bibitem[Carmona et~al.(2019)]{carmona2019}
Carmona, P., Climent, F., and Momparler, A.\ (2019).
\newblock Predicting failure in the U.S.\ banking sector: an extreme gradient
boosting approach.
\newblock \textit{International Review of Economics and Finance}, 61,
304--323.

\bibitem[Chawla et~al.(2002)]{chawla2002}
Chawla, N.\ V., Bowyer, K.\ W., Hall, L.\ O., and Kegelmeyer, W.\ P.\ (2002).
\newblock SMOTE: synthetic minority over-sampling technique.
\newblock \textit{Journal of Artificial Intelligence Research}, 16, 321--357.

\bibitem[Duan et~al.(2012)]{duan2012}
Duan, J.\ C., Sun, J., and Wang, T.\ (2012).
\newblock Multiperiod corporate default prediction --- a forward intensity
approach.
\newblock \textit{Journal of Econometrics}, 170(1), 191--209.

\bibitem[Fisher(1936)]{fisher1936}
Fisher, R.\ A.\ (1936).
\newblock The use of multiple measurements in taxonomic problems.
\newblock \textit{Annals of Eugenics}, 7(2), 179--188.

\bibitem[King and Zeng(2001)]{king2001}
King, G.\ and Zeng, L.\ (2001).
\newblock Logistic regression in rare events data.
\newblock \textit{Political Analysis}, 9(2), 137--163.

\bibitem[Mai et~al.(2019)]{mai2019}
Mai, F., Tian, S., Lee, C., and Ma, L.\ (2019).
\newblock Deep learning models for bankruptcy prediction using textual
disclosures.
\newblock \textit{European Journal of Operational Research}, 274(2),
743--758.

\bibitem[Min and Lee(2005)]{min2005}
Min, J.\ H.\ and Lee, Y.\ C.\ (2005).
\newblock Bankruptcy prediction using support vector machine with optimal
choice of kernel function parameters.
\newblock \textit{Expert Systems with Applications}, 28(4), 603--614.

\bibitem[Mossman et~al.(1998)]{mossman1998}
Mossman, C.\ E., Bell, G.\ G., Swartz, L.\ M., and Turtle, H.\ (1998).
\newblock An empirical comparison of bankruptcy models.
\newblock \textit{Financial Review}, 33(2), 35--54.

\bibitem[Niculescu-Mizil and Caruana(2005)]{niculescu2005}
Niculescu-Mizil, A.\ and Caruana, R.\ (2005).
\newblock Predicting good probabilities with supervised learning.
\newblock \textit{Proceedings of the 22nd International Conference on Machine
Learning}, 625--632.

\bibitem[Odom and Sharda(1990)]{odom1990}
Odom, M.\ D.\ and Sharda, R.\ (1990).
\newblock A neural network model for bankruptcy prediction.
\newblock \textit{Proceedings of the International Joint Conference on Neural
Networks}, 2, 163--168.

\bibitem[Ohlson(1980)]{ohlson1980}
Ohlson, J.\ A.\ (1980).
\newblock Financial ratios and the probabilistic prediction of bankruptcy.
\newblock \textit{Journal of Accounting Research}, 18(1), 109--131.

\bibitem[Pellegrino et~al.(2022)]{pellegrino2022}
Lombardo, G., Pellegrino, M., Adosoglou, G., Cagnoni, S., Pardalos, P.\ M.,
and Poggi, A.\ (2022).
\newblock Machine learning for bankruptcy prediction in the American stock
market: dataset and benchmarks.
\newblock \textit{Future Internet}, 14(8), 244.

\bibitem[Pellegrino et~al.(2024)]{pellegrino2024}
Pellegrino, M., Lombardo, G., Adosoglou, G., Cagnoni, S., Pardalos, P.\ M.,
and Poggi, A.\ (2024).
\newblock A multi-head LSTM architecture for bankruptcy prediction with
time-series accounting data.
\newblock \textit{Future Internet}, 16(3), 79.

\bibitem[Saito and Rehmsmeier(2015)]{saito2015}
Saito, T.\ and Rehmsmeier, M.\ (2015).
\newblock The precision-recall plot is more informative than the ROC plot
when evaluating binary classifiers on imbalanced datasets.
\newblock \textit{PLoS ONE}, 10(3), e0118432.

\bibitem[Saerens et~al.(2002)]{saerens2002}
Saerens, M., Latinne, P., and Decaestecker, C.\ (2002).
\newblock Adjusting the outputs of a classifier to new a priori probabilities:
a simple procedure.
\newblock \textit{Neural Computation}, 14(1), 21--41.

\bibitem[Shin et~al.(2005)]{shin2005}
Shin, K.\ S., Lee, T.\ S., and Kim, H.\ J.\ (2005).
\newblock An application of support vector machines in bankruptcy prediction
model.
\newblock \textit{Expert Systems with Applications}, 28(1), 127--135.

\bibitem[Shumway(2001)]{shumway2001}
Shumway, T.\ (2001).
\newblock Forecasting bankruptcy more accurately: a simple hazard model.
\newblock \textit{Journal of Business}, 74(1), 101--124.

\bibitem[Zięba et~al.(2016)]{zieba2016}
Zięba, M., Tomczak, S.\ K., and Tomczak, J.\ M.\ (2016).
\newblock Ensemble boosted trees with synthetic features generation in
application to bankruptcy prediction.
\newblock \textit{Expert Systems with Applications}, 58, 93--101.

\bibitem[Zmijewski(1984)]{zmijewski1984}
Zmijewski, M.\ E.\ (1984).
\newblock Methodological issues related to the estimation of financial
distress prediction models.
\newblock \textit{Journal of Accounting Research}, 22 (Supplement), 59--82.

\end{thebibliography}

\end{document}